\documentclass[11pt,a4paper]{article}

\usepackage[margin=1in]{geometry}
\usepackage{amsmath}
\usepackage{amssymb}
\usepackage{listings}
\usepackage{xcolor}
\usepackage{hyperref}
\usepackage{booktabs}

\lstset{
    language=C,
    basicstyle=\ttfamily\small,
    keywordstyle=\color{blue},
    commentstyle=\color{gray},
    stringstyle=\color{red},
    numbers=left,
    numberstyle=\tiny\color{gray},
    frame=single,
    breaklines=true,
    captionpos=b
}

\title{\textbf{TBITW ML Library}\\Implementation Details}
\author{dansyc11}
\date{}

\begin{document}

\maketitle

\noindent
\textbf{Last Updated:} May 11, 2026

\vspace{0.5cm}

\section{Introduction}

I built this library understand ML from first principle with just C

\subsection{Core Components}
\begin{itemize}
    \item Arena allocator for fast memory management
    \item Matrix operations for linear algebra
    \item Neural network with backpropagation
    \item Adam optimizer for training
\end{itemize}

\section{Memory Management}

\subsection{Arena Allocator}

Pre-allocate large block, bump pointer for O(1) allocation:

\begin{lstlisting}[caption={Arena Structure}]
typedef struct {
    u8 *memory;
    usize size;
    usize used;
} Arena;
\end{lstlisting}

Eliminates malloc/free overhead during training.

\begin{table}[h]
\centering
\begin{tabular}{lcc}
\toprule
Operation & malloc/free & Arena \\
\midrule
Single Allocation & O(log n) & O(1) \\
Batch Deallocation & O(n) & O(1) \\
\bottomrule
\end{tabular}
\caption{Allocator Performance}
\end{table}

\section{Neural Network}

\subsection{Architecture}

\begin{lstlisting}[caption={Network Structure}]
typedef struct {
    Mat *ws;     // Weights [L]
    Mat *bs;     // Biases [L]
    Mat *as;     // Activations [L+1]
    u32 count;
} NN;
\end{lstlisting}

\subsection{Forward Pass}

For layer $l$:
\begin{align}
    z^{(l)} &= a^{(l-1)} W^{(l)} + b^{(l)} \\
    a^{(l)} &= \sigma(z^{(l)})
\end{align}

Where sigmoid: $\sigma(x) = \frac{1}{1 + e^{-x}}$

\subsection{Cost Function}

Mean squared error:
\begin{equation}
    J = \frac{1}{N} \sum_{i=1}^{N} \sum_{j=1}^{n_L} (a_j^{(L)} - y_j)^2
\end{equation}

\section{Backpropagation}

Output layer gradient:
\begin{equation}
    \delta^{(L)} = 2(a^{(L)} - y) \odot \sigma'(z^{(L)})
\end{equation}

Hidden layer gradient:
\begin{equation}
    \delta^{(l)} = (\delta^{(l+1)} (W^{(l+1)})^T) \odot \sigma'(z^{(l)})
\end{equation}

Weight updates:
\begin{equation}
    \frac{\partial J}{\partial W^{(l)}} = (a^{(l-1)})^T \delta^{(l)}
\end{equation}

\section{Adam Optimizer}

Maintains momentum $m_t$ and adaptive learning rate $v_t$.

\textbf{Hyperparameters:} $\beta_1 = 0.9$, $\beta_2 = 0.999$, $\epsilon = 10^{-8}$, $\alpha = 0.001$

\textbf{Update rules:}
\begin{align}
    m_t &= \beta_1 m_{t-1} + (1-\beta_1) g_t \\
    v_t &= \beta_2 v_{t-1} + (1-\beta_2) g_t^2 \\
    \hat{m}_t &= \frac{m_t}{1-\beta_1^t}, \quad \hat{v}_t = \frac{v_t}{1-\beta_2^t} \\
    \theta_{t+1} &= \theta_t - \alpha \frac{\hat{m}_t}{\sqrt{\hat{v}_t} + \epsilon}
\end{align}

I implemented Adam to try fixing the XOR convergence problem. It didn't improve XOR performance significantly, but it did increase accuracy on both MNIST digits and Fashion-MNIST by 2-3\% compared to vanilla SGD.

\section{Training}

\subsection{Xavier Initialization}

\begin{equation}
    W^{(l)} \sim U\left[-\sqrt{\frac{6}{n_{l-1} + n_l}}, \sqrt{\frac{6}{n_{l-1} + n_l}}\right]
\end{equation}

Prevents gradient vanishing/explosion.

\subsection{Configuration}

\begin{table}[h]
\centering
\begin{tabular}{lc}
\toprule
Parameter & Value \\
\midrule
Learning Rate & 0.001 \\
Batch Size & 100 \\
Epochs & 20 \\
Architecture & 784-128-10 \\
\bottomrule
\end{tabular}
\end{table}

\section{Results}

\subsection{MNIST Performance}

\begin{table}[h]
\centering
\begin{tabular}{cccc}
\toprule
Epoch & Train Loss & Test Accuracy \\
\midrule
1 & 0.245 & 45.50\% \\
5 & 0.089 & 87.67\%  \\
10 & 0.032 & 90.68\%  \\
20 & 0.018 & 92.37\%  \\
\bottomrule
\end{tabular}
\caption{MNIST Training Progression}
\end{table}

\subsection{Fashion-MNIST}

Final accuracy: 92.16\% on 10 clothing categories.

\section{ASCII Visualization}

The prediction tools visualize 28x28 images in the terminal using ASCII characters. Each pixel value (0.0 to 1.0) is mapped to a character based on intensity:

\begin{lstlisting}[caption={ASCII Visualization}]
void print_digit(Mat image) {
    const char *chars = " .:oO@";
    for (u32 i = 0; i < 28; ++i) {
        for (u32 j = 0; j < 28; ++j) {
            f32 pixel = image.data[i * 28 + j];
            int idx = (int)(pixel * 5);
            printf("%c", chars[idx]);
        }
        printf("\n");
    }
}
\end{lstlisting}

Darker pixels (near 0.0) map to spaces, while brighter pixels (near 1.0) map to '@' symbols. This creates a recognizable visualization of handwritten digits and clothing items directly in the terminal.

\section{Model File Format}

Binary format:
\begin{enumerate}
    \item Magic: \texttt{0x54424954} ("TBIT")
    \item Layer count: \texttt{u32}
    \item For each layer: dimensions + weights + biases
\end{enumerate}

File size: $\sim$407 KB for 784-128-10 network.

\section{Future Improvements}

\begin{itemize}
    \item[$\square$] ReLU activation to improve gradient flow
    \item[$\square$] Convolutional layers to maximize MNIST accuracy (target: 98-99\%)
    \item[$\square$] Test on additional datasets beyond MNIST/Fashion-MNIST
\end{itemize}

\section{Appreciation}

\begin{itemize}
    \item \textbf{Tsoding} - ML in C series on YouTube. Core implementation techniques and educational approach taken from his tutorials on building neural networks from scratch.
    \item \textbf{Magicalbat} - Additional neural network implementation guidance and best practices.
\end{itemize}

\textit{Implementation written from scratch with educational guidance from the above resources.}

\end{document}
